% GENERATED FILE. Edit canonical lessons, then run npm run generate:books.
% canonical-source-hash: fnv1a64:a94f549282234835
% canonical-lessons: JA-C350-youjin-suru, JA-C350-yudan, JA-C350-sonaeru, JA-C350-taisaku, JA-C350-bousai

\chapter{Taking Care}
\label{ch:ja-taking-care}
\hlchaptermodality{350}

\begin{chapteropening}
\textbf{By the end of this chapter:} \emph{I can say what to do to stay safe if something happens.}

Complete the last lesson of chapter 350: I can say what to do to stay safe if something happens.
\end{chapteropening}

\section[\texorpdfstring{\ja{ようじんする} (yōjin suru)}{ようじんする (yōjin suru)}]{\ja{ようじんする} (yōjin suru) — to take precautions; to be on guard}

\label{lesson:JA-C350-youjin-suru}



\begin{quote}
Before the new one: say the Japanese for to miss a chance, then the Japanese for at great pains.
\end{quote}

\subsection*{You'll want to know: \ja{ようじんする}}
\textbf{\ja{ようじんする}} — \emph{yōjin suru} — \textquotedblleft{}to take precautions; to be on guard\textquotedblright{}.

\begin{quote}
\textbf{\ja{ばんに} \ja{でかけたら、ようじんして} \ja{ください。}} — \emph{ban ni dekaketara, yōjin shite kudasai} — \textquotedblleft{}If you go out in the evening, be on your guard.\textquotedblright{}
\end{quote}

Polite \textbf{\ja{ようじんします}}, \emph{yōjinshimasu}; a \ja{する} verb. Past (polite): \textbf{\ja{ちちは} \ja{ようじんしました。}}, \emph{chichi wa yōjin shimashita}.

\begin{cousinweb}
\textbf{\ja{ようじんする}} (\emph{yōjin suru}) is built from a Sino-Japanese noun and the native verb \emph{suru}, \textquotedblleft{}to do\textquotedblright{}. The noun, \emph{yōjin}, is two Chinese characters read with their on'yomi, readings borrowed from Chinese: \emph{yō}, \textquotedblleft{}utilize\textquotedblright{} (Middle Chinese \emph{yowngH}) and \emph{jin}, \textquotedblleft{}heart\textquotedblright{} (\emph{sim}) (KANJIDIC2; JMdict; Middle Chinese in Baxter's transcription).
\end{cousinweb}

\begin{grammarlens}[title={What you've built}]
One more word for saying what happens if.
\end{grammarlens}

\subsection*{Your turn}
\begin{itemize}
\raggedright
  \item \emph{Say it:} \emph{yōjin suru}
  \item \emph{Say it:} \emph{yōjin suru}, once more
  \item \emph{Say it:} \emph{ban ni dekaketara, yōjin shite kudasai} — If you go out in the evening, be on your guard.
  \item \emph{Recall:} say the Japanese for to miss a chance, then the Japanese for at great pains, then say \emph{yōjin suru} again
\end{itemize}

\begin{tcolorbox}[breakable,colback=teal!4,colframe=teal!35!black,title={Before you move on}]
What does \ja{ようじんする} mean? (\textquotedblleft{}To take precautions; to be on guard\textquotedblright{}.) Say it once more.
\end{tcolorbox}
\section[\texorpdfstring{\ja{ゆだん} (yudan)}{ゆだん (yudan)}]{\ja{ゆだん} (yudan) — carelessness; letting one's guard down}

\label{lesson:JA-C350-yudan}



\begin{quote}
Before the new one: say the Japanese for at great pains, then the Japanese for to take precautions.
\end{quote}

\subsection*{You'll want to know: \ja{ゆだん}}
\textbf{\ja{ゆだん}} — \emph{yudan} — \textquotedblleft{}carelessness; letting one's guard down\textquotedblright{}.

\begin{quote}
\textbf{\ja{ゆだんしたら、けがを} \ja{します。}} — \emph{yudan shitara, kega o shimasu} — \textquotedblleft{}If you let your guard down, you will get hurt.\textquotedblright{}
\end{quote}

\begin{cousinweb}
\textbf{\ja{ゆだん}} (\emph{yudan}) is Sino-Japanese (kango): two Chinese characters read with their on'yomi, readings borrowed from Chinese. They are \emph{yu}, \textquotedblleft{}oil\textquotedblright{} (Middle Chinese \emph{yuw}) and \emph{dan}, \textquotedblleft{}severance\textquotedblright{} (\emph{dwanX}) (KANJIDIC2; Middle Chinese in Baxter's transcription).
\end{cousinweb}

\begin{grammarlens}[title={What you've built}]
One more word for saying what happens if.
\end{grammarlens}

\subsection*{Your turn}
\begin{itemize}
\raggedright
  \item \emph{Say it:} \emph{yudan}
  \item \emph{Say it:} \emph{yudan}, once more
  \item \emph{Say it:} \emph{yudan shitara, kega o shimasu} — If you let your guard down, you will get hurt.
  \item \emph{Recall:} say the Japanese for at great pains, then the Japanese for to take precautions, then say \emph{yudan} again
\end{itemize}

\begin{tcolorbox}[breakable,colback=teal!4,colframe=teal!35!black,title={Before you move on}]
What does \ja{ゆだん} mean? (\textquotedblleft{}Carelessness; letting one's guard down\textquotedblright{}.) Say it once more.
\end{tcolorbox}
\section[\texorpdfstring{\ja{そなえる} (sonaeru)}{そなえる (sonaeru)}]{\ja{そなえる} (sonaeru) — to prepare for; to be ready for}

\label{lesson:JA-C350-sonaeru}



\begin{quote}
Before the new one: say the Japanese for to take precautions, then the Japanese for carelessness.
\end{quote}

\subsection*{You'll want to know: \ja{そなえる}}
\textbf{\ja{そなえる}} — \emph{sonaeru} — \textquotedblleft{}to prepare for; to be ready for\textquotedblright{}.

\begin{quote}
\textbf{\ja{たいふうが} \ja{きたら} \ja{こまるから、そなえた} \ja{ほうが} \ja{いいです。}} — \emph{taifū ga kitara komaru kara, sonaeta hō ga ii desu} — \textquotedblleft{}A typhoon would cause trouble, so we should be ready for one.\textquotedblright{}
\end{quote}

Polite \textbf{\ja{そなえます}}, \emph{sonaemasu}; an ichidan verb. Past (polite): \textbf{\ja{まちが} \ja{たいふうに} \ja{そなえました。}}, \emph{machi ga taifū ni sonaemashita}.

\begin{cousinweb}
\textbf{\ja{そなえる}} (\emph{sonaeru}) is a native Japanese word (wago), inherited rather than borrowed from Chinese, and its older origin is not settled. Its kanji was attached to it later as a kun'yomi, a native reading given to a Chinese character (JMdict, KANJIDIC2).
\end{cousinweb}

\begin{grammarlens}[title={What you've built}]
One more word for saying what happens if.
\end{grammarlens}

\subsection*{Your turn}
\begin{itemize}
\raggedright
  \item \emph{Say it:} \emph{sonaeru}
  \item \emph{Say it:} \emph{sonaeru}, once more
  \item \emph{Say it:} \emph{taifū ga kitara komaru kara, sonaeta hō ga ii desu} — A typhoon would cause trouble, so we should be ready for one.
  \item \emph{Recall:} say the Japanese for to take precautions, then the Japanese for carelessness, then say \emph{sonaeru} again
\end{itemize}

\begin{tcolorbox}[breakable,colback=teal!4,colframe=teal!35!black,title={Before you move on}]
What does \ja{そなえる} mean? (\textquotedblleft{}To prepare for; to be ready for\textquotedblright{}.) Say it once more.
\end{tcolorbox}
\section[\texorpdfstring{\ja{たいさく} (taisaku)}{たいさく (taisaku)}]{\ja{たいさく} (taisaku) — a countermeasure}

\label{lesson:JA-C350-taisaku}



\begin{quote}
Before the new one: say the Japanese for carelessness, then the Japanese for to prepare for.
\end{quote}

\subsection*{You'll want to know: \ja{たいさく}}
\textbf{\ja{たいさく}} — \emph{taisaku} — \textquotedblleft{}a countermeasure\textquotedblright{}.

\begin{quote}
\textbf{\ja{あめが} \ja{ふったら、どんな} \ja{たいさくが} \ja{ありますか。}} — \emph{ame ga futtara, donna taisaku ga arimasu ka} — \textquotedblleft{}If it rains, what can we do about it?\textquotedblright{}
\end{quote}

\begin{cousinweb}
\textbf{\ja{たいさく}} (\emph{taisaku}) is Sino-Japanese (kango): two Chinese characters read with their on'yomi, readings borrowed from Chinese. They are \emph{tai}, \textquotedblleft{}opposite\textquotedblright{} (Middle Chinese \emph{twojH}) and \emph{saku}, \textquotedblleft{}scheme\textquotedblright{} (\emph{tsrheak}) (KANJIDIC2; Middle Chinese in Baxter's transcription).
\end{cousinweb}

\begin{grammarlens}[title={What you've built}]
One more word for saying what happens if.
\end{grammarlens}

\subsection*{Your turn}
\begin{itemize}
\raggedright
  \item \emph{Say it:} \emph{taisaku}
  \item \emph{Say it:} \emph{taisaku}, once more
  \item \emph{Say it:} \emph{ame ga futtara, donna taisaku ga arimasu ka} — If it rains, what can we do about it?
  \item \emph{Recall:} say the Japanese for carelessness, then the Japanese for to prepare for, then say \emph{taisaku} again
\end{itemize}

\begin{tcolorbox}[breakable,colback=teal!4,colframe=teal!35!black,title={Before you move on}]
What does \ja{たいさく} mean? (\textquotedblleft{}A countermeasure\textquotedblright{}.) Say it once more.
\end{tcolorbox}
\section[\texorpdfstring{\ja{ぼうさい} (bōsai)}{ぼうさい (bōsai)}]{\ja{ぼうさい} (bōsai) — disaster preparedness}

\label{lesson:JA-C350-bousai}



\begin{quote}
Before the new one: say the Japanese for to prepare for, then the Japanese for a countermeasure.
\end{quote}

\subsection*{You'll want to know: \ja{ぼうさい}}
\textbf{\ja{ぼうさい}} — \emph{bōsai} — \textquotedblleft{}disaster preparedness\textquotedblright{}.

\begin{quote}
\textbf{\ja{もし} \ja{じしんが} \ja{きたら、ぼうさいの} \ja{きそくに} \ja{したがって} \ja{ください。}} — \emph{moshi jishin ga kitara, bōsai no kisoku ni shitagatte kudasai} — \textquotedblleft{}If an earthquake comes, follow the disaster rules.\textquotedblright{}
\end{quote}

\begin{cousinweb}
\textbf{\ja{ぼうさい}} (\emph{bōsai}) is Sino-Japanese (kango): two Chinese characters read with their on'yomi, readings borrowed from Chinese. They are \emph{bō}, \textquotedblleft{}ward off\textquotedblright{} (Middle Chinese \emph{bjang}) and \emph{sai}, \textquotedblleft{}disaster\textquotedblright{} (\emph{tsoj}) (KANJIDIC2; Middle Chinese in Baxter's transcription).
\end{cousinweb}

\begin{grammarlens}[title={What you've built}]
One more word for saying what happens if.
\end{grammarlens}

\subsection*{Your turn}
\begin{itemize}
\raggedright
  \item \emph{Say it:} \emph{bōsai}
  \item \emph{Say it:} \emph{bōsai}, once more
  \item \emph{Say it:} \emph{moshi jishin ga kitara, bōsai no kisoku ni shitagatte kudasai} — If an earthquake comes, follow the disaster rules.
  \item \emph{Recall:} say the Japanese for to prepare for, then the Japanese for a countermeasure, then say \emph{bōsai} again
\end{itemize}

\begin{tcolorbox}[breakable,colback=teal!4,colframe=teal!35!black,title={Before you move on}]
What does \ja{ぼうさい} mean? (\textquotedblleft{}Disaster preparedness\textquotedblright{}.) Say it once more.
\end{tcolorbox}
